%======================================================================
%  Workload concepts and analytical model
%======================================================================

\section{A performance model for LLM inference}
\label{sec:inference-model}

\subsection{Prefill and autoregressive decode}

Inference for a decoder-only transformer has two phases
~\cite{na2024llmcpu}.  During
\emph{prefill}, the model processes the prompt and constructs the
initial key--value (KV) cache.  Tokens from the prompt can be processed
in parallel, so the linear layers often form comparatively large matrix
multiplications.  During \emph{decode}, the model generates one new
token per sequence at each iteration.  Each iteration requires the
relevant weights and KV state, then appends the key and value for the new
token.  Dense full-context attention reads the prior cached KV; sliding
window or sparse attention can reduce that history.
\figref{fig:prefill-decode} shows the two phases and the metric that
characterises each.

\begin{figure}[htbp]
\centering
\begin{tikzpicture}[
  font=\small,
  io/.style={draw=black!50, fill=white, rounded corners=2pt,
             minimum width=2.2cm, minimum height=0.9cm, align=center},
  phase/.style={draw=black!65, fill=neutralcol, rounded corners=3pt,
                minimum width=4.0cm, minimum height=1.55cm,
                text width=3.7cm, align=center},
  kv/.style={draw=hbmcol!85!black, fill=hbmcol!16, rounded corners=2pt,
             minimum width=4.0cm, minimum height=0.8cm, align=center},
  flow/.style={-{Stealth}, thick},
  note/.style={font=\footnotesize, align=center}
]
  \node[io] (prompt) at (-6.0,0) {prompt};
  \node[phase] (prefill) at (-2.5,0)
    {\textbf{Prefill}\\process prompt in parallel\\build initial KV};
  \node[phase] (decode) at (2.6,0)
    {\textbf{Decode loop}\\access weights and KV\\
     generate one token\\per sequence};
  \node[io] (output) at (6.3,0) {token stream};
  \node[kv] (cache) at (2.6,-2.1)
    {KV cache: read required history, append new state};

  \draw[flow] (prompt) -- (prefill);
  \draw[flow] (prefill) -- (decode);
  \node[note, black!75] at (0.05,1.45) {first token};
  \draw[flow] (decode) -- (output);
  \draw[flow, {Stealth}-{Stealth}] (decode) -- (cache);
  \draw[flow] (decode.45) .. controls +(0.8,0.8) and +(-0.8,0.8)
       .. (decode.135) node[midway, above, note]{next iteration};

  \node[note, below=0.2cm of prefill]
    {principal latency metric: TTFT};
  \node[note, below=0.65cm of cache]
    {principal latency metric: inter-token latency};
\end{tikzpicture}
\caption{Dataflow of prompt prefill and autoregressive decode.  The
frequent description of prefill as compute-bound and decode as
bandwidth-bound is a useful first-order model, not a universal property:
batch, context, precision and kernel implementation can change the
regime.  Author's synthesis.}
\label{fig:prefill-decode}
\end{figure}

This phase distinction is essential when evaluating HBM.  A single
end-to-end throughput number can hide a faster decode phase behind a
long prompt, or hide a prefill regression behind a long generation.
\textbf{TTFT and decode latency must therefore be reported
separately.}  Prefill
usually offers matrix--matrix parallelism, whereas low-batch decode is
closer to matrix--vector execution and repeatedly streams weights.
Larger batches increase weight reuse and can move decode toward
matrix-engine throughput; long contexts simultaneously increase KV
traffic.


\subsection{Working-set capacity}
\label{sec:working-set}

For a model with \(P\) parameters and a stored weight precision of
\(q_w\) bits, \eqnref{eq:weights} gives the first-order weight
footprint:

\begin{equation}
\label{eq:weights}
  M_W = P\frac{q_w}{8} + M_{\mathrm{format}},
\end{equation}
\eqcaption{Resident weight footprint.  The second term is the one that surprises people: it collects scales, zero-points, alignment and any prepacked copy the runtime builds, so a GGUF file of a given size does not occupy that size in memory.}

where \(M_{\mathrm{format}}\) includes scales, zero-points, alignment,
and any prepacked or duplicated representation created by the runtime.
The model file size is \emph{not} the resident weight footprint, and
treating the two as interchangeable is the commonest way to misjudge
whether a model fits a tier.

For \(B\) concurrent sequences, cached length \(S\), \(L\) transformer
layers, \(H_{\mathrm{KV}}\) KV heads, head dimension \(d_h\), and
\(q_{\mathrm{KV}}\) bits per cached element, \eqnref{eq:kv} gives the
KV footprint:

\begin{equation}
\label{eq:kv}
  M_{\mathrm{KV}} =
  2 B S L H_{\mathrm{KV}} d_h \frac{q_{\mathrm{KV}}}{8}.
\end{equation}
\eqcaption{Key--value cache footprint.  It grows linearly in both context length \(S\) and batch \(B\), which is why long prompts and concurrency are the two ways to leave a memory tier.}

The factor two accounts for keys and values.  Grouped-query and
multi-query attention reduce the KV footprint because
\(H_{\mathrm{KV}}\) is smaller than the number of query heads.  The
complete resident set is

\begin{equation}
\label{eq:live}
  M_{\mathrm{live}} =
  M_W + M_{\mathrm{KV}} + M_{\mathrm{activations}}
  + M_{\mathrm{workspace}} + M_{\mathrm{runtime}}.
\end{equation}
\eqcaption{The live set.  This, not the model file, is the quantity that must be compared against usable local HBM.}

This arithmetic identifies where capacity boundaries may occur, but the
boundary must use the OS-exposed NUMA capacity on the allocated node and
a measured safety margin.  The advertised 64\,GB per socket is not a
usable-allocation guarantee: system use, allocator fragmentation and
workspace consume part of it.  A node may expose roughly 128\,GB across
two sockets, but that capacity is not uniformly local; using it requires
a placement strategy that distributes data and execution deliberately.

Quantization changes both sides of the experiment.  Fewer bits reduce
weight capacity and memory traffic, potentially reducing the benefit of
HBM for a model that already fits.  At the same time, quantization can
move a larger model below the HBM capacity boundary and remove DDR or
remote traffic altogether.  Packing, dequantization and kernel coverage
prevent performance from scaling exactly with bit width.  Weight and KV
precision must therefore be reported separately.

\subsection{Arithmetic intensity and the Roofline bound}
\label{sec:roofline}

The Roofline model, stated as \eqnref{eq:roofline}, relates arithmetic
intensity \(I=F/D\), measured in FLOP/byte, to compute throughput and
memory bandwidth~\cite{williams2009roofline}:

\begin{equation}
\label{eq:roofline}
  P_{\mathrm{attainable}}
  \leq
  \min\!\left(P_{\mathrm{compute}},
              I B_{\mathrm{effective}}\right).
\end{equation}
\eqcaption{The Roofline bound.  Performance is capped by whichever runs out first, arithmetic or bandwidth; changing memory tier moves only the second term.}

\eqnref{eq:step} is the corresponding lower-bound model for a single
inference step:

\begin{equation}
\label{eq:step}
  T_{\mathrm{step}}
  \gtrsim
  \max\!\left(
      \frac{F_{\mathrm{step}}}{P_{\mathrm{effective}}},
      \frac{D_{\mathrm{step}}}{B_{\mathrm{effective}}}
  \right)
  + T_{\mathrm{communication}} + T_{\mathrm{runtime}}.
\end{equation}
\eqcaption{Lower bound on one inference step.  The two trailing terms are the reason a measured speedup can fall well short of a bandwidth ratio even when the step really is memory-bound.}

The effective quantities must be measured for the selected kernel,
placement and core count.  Processor peak FLOP/s and theoretical pin
bandwidth are ceilings, not values to substitute uncritically.

For a dense linear layer during decode, if a stored weight is used once
for each of \(B\) sequences in a batch, \eqnref{eq:intensity} is a
useful first-order approximation:

\begin{equation}
\label{eq:intensity}
  I_{\mathrm{weights}} \approx \frac{2B}{b_w}
  \quad \text{FLOP/byte},
\end{equation}
\eqcaption{Arithmetic intensity of weight streaming in decode.  Batch size is in the numerator and bytes per weight in the denominator, so batching moves right along the Roofline and quantization moves right as well, both away from the region where a memory tier decides the outcome.}

where \(b_w\) is the stored bytes per weight and a multiply--add is
counted as two FLOPs.  The approximation assumes that each fetched dense
weight is reused across the \(B\) sequences.  At batch one, BF16 weight
streaming is therefore close to one FLOP/byte before KV,
activation and runtime traffic are included.  Batching reuses weights
and raises arithmetic intensity, which explains why HBM is expected to
matter most for low-batch decode and why AMX can become more important
as the matrix dimensions grow.

\figref{fig:roofline} makes the consequence visible.  Placing both
memory tiers on the same Roofline separates the region where a placement
decision can matter from the region where it cannot: to the left of the
ridge points the vertical distance between the two roofs is the full
bandwidth ratio, while to the right both roofs lie above the compute
ceiling, so no placement policy changes the bound.  Low-batch decode
sits in the first region and prefill in the second.  That is the
argument of this thesis expressed as a picture, and it is also why the
two phases have to be measured separately rather than folded into a
single throughput figure.

\begin{figure}[htbp]
\centering
\begin{tikzpicture}[
  font=\small, x=1.22cm, y=1.0cm,
  roof/.style={very thick},
  point/.style={circle, fill=black!80, inner sep=1.6pt},
  lbl/.style={font=\footnotesize, align=center}
]
  \draw[->, black!70] (0,0) -- (9.35,0)
    node[right, font=\scriptsize, align=center]
      {FLOP/\\byte};
  \draw[->, black!70] (0,0) -- (0,4.05)
    node[above, font=\footnotesize] {attainable GFLOP/s};

  \foreach \X/\lab in {0/0.5, 0.818/1, 2.455/4, 4.091/16,
                       5.727/64, 7.364/256, 9.0/1024}
    \draw[black!55] (\X,0) -- (\X,-0.09) node[below, font=\scriptsize] {\lab};

  \foreach \Y/\lab in {0/{10^2}, 1.08/{10^3}, 2.16/{10^4}, 3.24/{10^5}}
    \draw[black!55] (0,\Y) -- (-0.10,\Y)
      node[left, font=\scriptsize] {$\lab$};

  \draw[black!12] (0,1.08) -- (9,1.08);
  \draw[black!12] (0,2.16) -- (9,2.16);
  \draw[black!12] (0,3.24) -- (9,3.24);

  \fill[black!5] (1.6,0) rectangle (3.3,3.28);
  \draw[black!30, dashed] (1.6,0) -- (1.6,3.28);
  \draw[black!30, dashed] (3.3,0) -- (3.3,3.28);

  \draw[dashed, thick, black!60] (0,3.28) -- (9,3.28);
  \node[lbl, black!65, anchor=south east] at (9,3.36)
    {AMX BF16 peak (specification ceiling)};

  \draw[roof, hbmcol!85!black] (0,0.673) -- (6.56,3.28) -- (9,3.28);
  \node[lbl, hbmcol!65!black, rotate=30, anchor=south west]
    at (3.55,1.72) {HBM $\approx$ 840\,GB/s};

  \draw[roof, ddrcol!85!black] (0,0.095) -- (8.02,3.28) -- (9,3.28);
  \node[lbl, ddrcol!75!black, rotate=26, anchor=north west]
    at (4.35,1.62) {DDR5 $\approx$ 245\,GB/s};

  \node[point, hbmcol!70!black] at (2.455,1.65) {};
  \node[point, ddrcol!75!black] at (2.455,1.07) {};
  \draw[{Stealth}-{Stealth}, thick, black!75] (2.455,1.07) -- (2.455,1.65);
  \node[font=\scriptsize, align=center, black!80, fill=white,
        inner sep=1.5pt] at (2.455,0.58) {$\approx$3.4$\times$ apart};

  \node[lbl, black!75] at (2.45,-0.80)
    {batch-1 decode,\\4-bit weights};
  \node[lbl, black!75] at (5.85,-0.80)
    {batched decode,\\longer context};
  \node[lbl, black!75] at (8.45,-0.80) {prefill\\(large GEMM)};

  \node[point] at (5.727,3.28) {};
  \node[point] at (8.30,3.28) {};
  \draw[black!45, -{Stealth}] (5.85,-0.40) -- (5.75,3.15);
  \draw[black!45, -{Stealth}] (8.45,-0.40) -- (8.32,3.15);

\end{tikzpicture}
\caption{Both memory tiers of one Xeon Max socket on a single Roofline.
The bandwidth figures are full-socket Flat-mode STREAM Triad values
measured on a Xeon Max 9470C~\cite{ibeid2025aurora} and are used only to
set the slopes; the compute ceiling is a specification figure rather
than a measurement, and neither has yet been established for CRESCO8.
The shaded band marks the operating region this thesis investigates:
left of the ridge points the vertical distance between the roofs is the
full bandwidth ratio, while past them both roofs coincide with the
compute ceiling and placement cannot move the bound.  Author's
synthesis.}
\label{fig:roofline}
\end{figure}


This model also clarifies why bandwidth alone is insufficient:

\begin{itemize}
  \item \textbf{Bandwidth} sets the rate for many concurrent,
        sufficiently regular memory requests.
  \item \textbf{Latency} sets the cost of a dependent or lightly loaded
        access path.  HBM can have higher unloaded latency than DDR5
        while sustaining much greater bandwidth under load
        ~\cite{ibeid2025aurora}.
  \item \textbf{Capacity} determines whether the critical working set
        remains local.  Crossing a 64\,GB socket boundary can introduce
        DDR5 or UPI traffic and cause a discontinuity that a simple
        single-tier Roofline does not predict.
\end{itemize}

For concurrent HBM and DDR5 traffic, a no-contention service-time
estimate can range from perfect overlap to serialization of the
per-tier transfers.  It is not a general bound: shared uncore resources,
cache conflicts, UPI traffic and memory-controller interference can make
the measured time worse than the serialized estimate.  Simultaneous
HBM--DDR traffic must therefore be characterized before a tiered
placement policy is modelled.

\subsection{Placement is part of the algorithm}

In Flat mode the operating system exposes memory tiers through NUMA,
but it does not know which model objects are bandwidth-critical.
Several mechanisms can therefore produce the wrong placement:
allocator first touch, worker migration, memory-mapped model pages,
automatic NUMA balancing, and a runtime that creates prepacked copies
after the initial allocation.  A command-line memory policy is evidence
of intent, \textbf{not evidence of residency}.

The relevant experimental object is consequently not ``the model in
HBM'' but a verified mapping between:

\begin{itemize}
  \item worker threads and physical cores;
  \item weight, KV, activation and workspace pages and their NUMA nodes;
  \item local or remote memory traffic observed during each inference
        phase; and
  \item the time interval used for performance and energy measurement.
\end{itemize}

This mapping is also the basis for testing topology-aware execution.
Independent replicas can use one local socket each; a single model can
be sharded across sockets or SNC domains; and selected objects can be
placed in different memory tiers.  These designs answer different
questions and should not be grouped under a generic ``two-socket run''.

\subsection{What we measure, and what each number means}
\label{sec:metrics}

A run of the experiment produces three kinds of number, and the most
common way to make a memory-placement study unreadable is to mix them.
A \textbf{result} is a quantity the thesis is willing to defend as an
outcome.  An \textbf{explanation} is a quantity that says \emph{why}
the outcome came out that way, and carries no weight on its own.  A
\textbf{validity check} is a quantity that decides whether the run
counts at all.  The tables below keep them apart, and every reported
comparison is expected to carry all three.

\subsubsection{Numbers that are results}

{\small\tabsetup
\begin{xltabular}{\textwidth}{@{}P{2.6cm}P{4.1cm}Y@{}}
\caption{Performance metrics: the quantities a claim may be made about.
A metric whose movement cannot be interpreted is not worth collecting,
so the right-hand column is the one to read first.}
\label{tab:metrics-results}\\
\toprule
\headrow
\textbf{Metric} & \textbf{What it is} & \textbf{Why it is here, and how to read it} \\
\midrule
\endfirsthead
\multicolumn{3}{@{}l}{\footnotesize\itshape Table~\ref{tab:metrics-results}, continued}\\[2pt]
\toprule
\headrow
\textbf{Metric} & \textbf{What it is} & \textbf{Why it is here, and how to read it} \\
\midrule
\endhead
\midrule
\multicolumn{3}{r@{}}{\footnotesize\itshape continued on the next page}\\
\endfoot
\bottomrule
\endlastfoot
\textbf{TTFT}\newline{\footnotesize time to first token}
& Elapsed time from an accepted request to the first emitted token.
& The visible cost of \emph{prefill}, and the only phase metric a long
  generation cannot hide.  \emph{Reading it:} scales with prompt
  length.  Because prefill sits right of the ridge point in
  \figref{fig:roofline}, a \textbf{large} tier effect here would be
  surprising and would need explaining. \\
\addlinespace[2pt]
\textbf{ITL}\newline{\footnotesize inter-token latency}
& The gap between consecutive emitted tokens, recorded per token.
& The decode metric, and where \(H_{1a}\) predicts the HBM effect will
  appear.  \emph{Reading it:} report the \textbf{distribution}
  (median, p95, p99), never a bare mean.  A stable median with a heavy
  tail points to allocation or scheduling interference, not to memory
  bandwidth. \\
\addlinespace[2pt]
\textbf{TPOT}\newline{\footnotesize time per output token}
& Mean decode time,
  \((T_{\mathrm{E2E}} - T_{\mathrm{TTFT}})/(O-1)\).
& A single summary of decode, and the metric the closest competing work
  reports~\cite{zhang2026cacheresident}.  \emph{Reading it:} good for
  comparison against published numbers, but it hides variance by
  construction, so it never appears without the ITL distribution beside
  it. \\
\addlinespace[2pt]
\textbf{Decode throughput}
& Output tokens per second, per sequence and aggregated over the batch.
& The form most readers expect and most CPU-inference papers publish.
  \emph{Reading it:} the reciprocal of TPOT, carrying no information
  TPOT does not; reported for comparability rather than for analysis. \\
\addlinespace[2pt]
\textbf{Prefill throughput}
& Prompt tokens per second, \(S/T_{\mathrm{TTFT}}\).
& Separates processing rate from prompt length, keeping runs with
  different prompts comparable.  \emph{Reading it:} roughly flat in
  prompt length once the prompt is long enough to fill the matrix
  units; a fall-off indicates a capacity or kernel effect. \\
\addlinespace[2pt]
\textbf{Tier ratio}\newline{\footnotesize \(R_{\text{phase}}\)}
& Per-phase ratio of matched DDR and HBM timings, \eqnref{eq:ratio}.
& \textbf{The headline result of the thesis.}  Everything else exists
  to make this number interpretable.  \emph{Reading it:}
  \(R_{\text{decode}} > R_{\text{prefill}}\) supports \(H_{1a}\); a
  \(R_{\text{decode}}\) below the STREAM bandwidth ratio supports
  \(H_{1b}\).  A value near \(1.0\) is \emph{a real finding}, not a
  failed experiment. \\
\end{xltabular}}


The tier ratio is defined on matched runs that differ in nothing but
placement:

\begin{equation}
\label{eq:ratio}
  R_{\text{phase}} =
  \frac{T_{\text{phase}}^{\text{DDR}}}{T_{\text{phase}}^{\text{HBM}}},
  \qquad \text{phase} \in \{\text{prefill}, \text{decode}\}.
\end{equation}
\eqcaption{Per-phase speedup from placing the live set in local HBM
rather than local DDR5.  Values above one mean HBM is faster.  Stated
per phase because a single end-to-end ratio would average the two
regimes of \figref{fig:roofline} into a number that describes neither.}

\subsubsection{Numbers that explain a result}

\tabref{tab:metrics-explain} lists the measurements that give a
result its mechanism.  The second row is the one to reach for first
when an effect turns out smaller than expected.


{\small\tabsetup
\begin{xltabular}{\textwidth}{@{}P{2.6cm}P{4.1cm}Y@{}}
\caption{Explanatory measurements.  None of these is a performance
claim.  Their job is to say why the numbers in
\tabref{tab:metrics-results} came out as they did, and, when the effect
is small, to separate \emph{the tier does not matter} from \emph{the
program never used the tier}.}
\label{tab:metrics-explain}\\
\toprule
\headrow
\textbf{Measurement} & \textbf{What it is} & \textbf{Why it is here, and how to read it} \\
\midrule
\endfirsthead
\multicolumn{3}{@{}l}{\footnotesize\itshape Table~\ref{tab:metrics-explain}, continued}\\[2pt]
\toprule
\headrow
\textbf{Measurement} & \textbf{What it is} & \textbf{Why it is here, and how to read it} \\
\midrule
\endhead
\midrule
\multicolumn{3}{r@{}}{\footnotesize\itshape continued on the next page}\\
\endfoot
\bottomrule
\endlastfoot
\textbf{Achieved bandwidth}
& Bytes per second actually moved to and from each tier during a phase,
  taken from memory-controller counters.
& \textbf{The most important explanatory number in the study};
  without it a null result cannot be diagnosed at all.
  \emph{Reading it:} compare against the STREAM figure measured on the
  same node, never against a published one. \\
\addlinespace[2pt]
\textbf{Bandwidth utilisation}
& Achieved bandwidth as a fraction of measured STREAM,
  \eqnref{eq:utilisation}.
& Turns the argument of \secref{sec:littles-law} into a measurement.
  \emph{Reading it:} a high tier ratio with high utilisation means the
  tier was the limit; a low tier ratio with \textbf{low} utilisation
  means \emph{concurrency} was the limit and the tier never got a
  chance. \\
\addlinespace[2pt]
\textbf{Remote traffic}
& Bytes crossing UPI to the other socket.
& The primary comparison is defined on one socket, so remote traffic
  means it was not.  \emph{Reading it:} expected near zero; anything
  else is a placement failure, and the run is re-examined rather than
  reported. \\
\addlinespace[2pt]
\textbf{Peak resident set}
& Maximum resident memory, decomposed into weights, KV state and
  workspace per \eqnref{eq:live}.
& Locates the capacity boundary the second research question is about.
  \emph{Reading it:} compared against \emph{usable} local HBM on the
  allocated node, which is smaller than the nominal 64\,GB. \\
\addlinespace[2pt]
\textbf{Frequency and power}
& Sustained socket clock, package power and throttling state over the
  measured interval.
& Guards against reporting a thermal effect as a placement effect.
  \emph{Reading it:} must be equivalent across the two placements being
  compared; if it is not, the comparison is confounded. \\
\end{xltabular}}


\begin{equation}
\label{eq:utilisation}
  U_{\text{tier}} =
  \frac{B_{\text{achieved}}}{B_{\text{STREAM}}} .
\end{equation}
\eqcaption{Bandwidth utilisation: the fraction of a tier's measured
sustainable bandwidth that the workload actually drew.  This is the
number that decides how a disappointing tier ratio should be read.  A
decode phase running at \(U_{\mathrm{HBM}} \approx 0.3\) has not shown
that HBM fails to help; it has shown that the workload asked HBM for
nothing DDR5 could not also have supplied.}

\subsubsection{Numbers that decide whether a run counts}

The checks in \tabref{tab:metrics-validity} come first in the analysis,
not last.  A run that fails any of them is discarded before its timings
are looked at, because a placement study whose placement was not
verified is not evidence of anything.

{\small\tabsetup
\begin{xltabular}{\textwidth}{@{}P{2.55cm}YP{4.3cm}@{}}
\caption{Validity checks applied to every run before its results are
considered.  These are pass/fail, not quantities to be reported as
findings.}
\label{tab:metrics-validity}\\
\toprule
\headrow
\textbf{Check} & \textbf{What it establishes} &
\textbf{Passing condition} \\
\midrule
\endfirsthead
\multicolumn{3}{@{}l}{\footnotesize\itshape Table~\ref{tab:metrics-validity}, continued}\\[2pt]
\toprule
\headrow
\textbf{Check} & \textbf{What it establishes} &
\textbf{Passing condition} \\
\midrule
\endhead
\midrule
\multicolumn{3}{r@{}}{\footnotesize\itshape continued on the next page}\\
\endfoot
\bottomrule
\endlastfoot
\textbf{Page residency}
& That the live set is where the command line said it should be.  A
  memory policy states an intention; only a page count states an
  outcome.
& The overwhelming majority of the process's resident pages sit on the
  intended NUMA node, with the threshold fixed before the campaign and
  the shortfall reported. \\
\addlinespace[1pt]
\textbf{Thread binding}
& That the cores executing the work are local to the memory holding it.
& Observed CPU affinity matches the requested binding for every worker
  thread. \\
\addlinespace[1pt]
\textbf{Output equivalence}
& That placement changed the speed and nothing else.
& For a fixed seed and identical inputs, the generated token sequence
  is identical across tiers.  A divergence indicates a defect, not a
  finding. \\
\addlinespace[1pt]
\textbf{Repetition dispersion}
& That an observed difference is larger than the noise it sits in.
& Coefficient of variation across repetitions small enough that the
  tier effect exceeds it; repetitions, warm-up and discarded runs all
  reported. \\
\addlinespace[1pt]
\textbf{Thermal state}
& That the node was in a comparable operating condition for both arms.
& No throttling events, and equivalent sustained frequency across the
  compared placements. \\
\end{xltabular}}


\subsubsection{Where each timing starts and stops}

The two phase metrics compose, as \eqnref{eq:e2e} shows, into the
latency a user actually experiences over \(O\) generated tokens in one
synchronous request:

\begin{equation}
\label{eq:e2e}
  T_{\mathrm{E2E}}
  =
  T_{\mathrm{TTFT}} +
  \sum_{j=2}^{O} T_{\mathrm{ITL},j}.
\end{equation}
\eqcaption{End-to-end request latency as the sum of its two phases.
Written this way to make the risk explicit: a single end-to-end figure
can absorb a decode improvement into a long prompt, or a prefill
regression into a long generation, which is why this thesis reports the
two terms rather than their sum.}

A latency figure is meaningless without its boundary, and boundaries
differ enough between published studies that comparing numbers across
papers is often comparing definitions~\cite{agrawal2025evaluating}.
This thesis fixes them as follows.  Model loading, tokenizer
construction and the first warm-up request lie \emph{outside} every
reported timing.  TTFT starts when the prompt has been tokenized and
submitted, and stops when the first token is available.  ITL is
measured between token emissions and excludes detokenization of the
final output.  Every reported figure comes from a warmed process, and
the number of discarded warm-up iterations is recorded with the run.

Two further distinctions are worth stating because they are routinely
elided.  Static batch throughput and online serving goodput are
different constructs: the second depends on an arrival process, a
scheduler and a declared latency objective, none of which are part of
this study.  And hardware counters are explanatory measurements rather
than performance metrics.  Their collection overhead, multiplexing
behaviour and measurement scope are documented alongside the values, or
the values are not used.

\subsection{Predictions to test}

The model yields three bounded predictions:

\begin{enumerate}
  \item local HBM should benefit low-batch decode more than compute-dense
        prefill, but the speedup need not equal the STREAM bandwidth
        ratio;
  \item performance should change when the verified live set crosses a
        local HBM boundary, with the size and direction of the change
        depending on placement and memory mode; and
  \item batching and lower-bit weights should reduce weight-bandwidth
        pressure, while long contexts can reintroduce pressure through
        the KV cache.
\end{enumerate}

These statements are hypotheses.  The literature reviewed next
motivates them but does not establish them for CRESCO8.
